\par\Needspace{7\baselineskip}\bigskip
\begin{center}
\pdfbookmark[2]{Gauss an Olbers, Göttingen, 12. August 1811}{pallas-brief-18110812}
\textsc{Gauss} an \textsc{Olbers}, Göttingen, 12. August 1811.
\end{center}
\noindent\hspace*{\parindent}. . . . .\quad Ich habe Ihnen . . . .\ nicht eher geschrieben, um doch noch Einiges über Ihre Pallas hinzusetzen zu können. Einiges, was Sie in einem der nächsten Stücke unserer Gel. Anz. finden werden\HistoricalNoteEight{Band VI, S. 327.]} (ich glaube vom 17. Aug.), übergehe ich, um nur etwas von den allgemeinen Perturbationsresultaten zu erwähnen. Ich habe die Störungen der Neigung und der Länge der Knotenlinie (als das am wenigsten mühsame) in der ersten Rechnung vollendet, auch erstere an den 7 Oppositionen (für welche meine frühern Arbeiten die Elemente geben) geprüft. Anfangs schlechte Übereinstimmung, bis ich fand, dass 2 Gleichungen das unrechte Zeichen hatten. Jetzt stimmen sie zu meiner Zufriedenheit überein. Mit der Prüfung des \(\AscendingNode\) werde ich auch in einigen Tagen fertig werden, dann will ich die Excentricität vornehmen, eine Arbeit von einem Monat etwa. Um Ihnen doch eine Idee von den Resultaten zu geben, füge ich die 40 Störungsgleichungen für die Neigung bei; ich habe bloss beibehalten, was über Eine Secunde ging; auch können vielleicht noch einige kleine Aequatiunculae, die von der 9 fachen und 10 fachen Jupiterslänge abhangen, dazu kommen\HistoricalNoteDouble{Hier folgen die S. 518 dieses Bandes gegebenen Glieder von \(\delta i\).]}.\quad . . . . .
\par
Die 80 Gleichungen für Inclination und \(\AscendingNode\) liessen sich in 40 für die Breite zusammenziehen, ich glaube aber nicht, dass etwas gewonnen wird, denn wenn man die Elemente selbst stören lässt, so kann man ohne Bedenken einerlei gestörte Elemente als mehrere Monate gültig ansehen, und braucht also alle Jahr nur einmal für 6 Elemente die Störungen zu berechnen (vielleicht zusammen etwa 300—400 Gleichungen); dahingegen, wenn man bei den Elementen bloss Secularänderungen anbringt und die periodischen bei Breite, Länge und Radius Vector (zusammen vielleicht gegen 200), diese in einem Jahre doch wohl wenigstens für 6 verschiedene Örter berechnet werden müssten, um interpoliren zu können. Doch kann man dies in der \SourceLeaf{427}Folge machen, wie man will, wenn nur erst alle Störungen \textls[180]{in irgend einer Form} da sind.
\par
